\documentclass[preview]{standalone}

\usepackage{amsmath}
\usepackage{amssymb}
\usepackage{stellar}
\usepackage{definitions}
\usepackage{bettelini}

\begin{document}

\id{probability-expectation}
\genpage

\section{Expected value}

\begin{snippetdefinition}{expected-value-discrete-random-variable-definition}{Expected value of a discrete random variable}
    Let \((\Omega,\powerset(\Omega),\mathbb P)\) be a
    \countableprobabilityspace, and let
    \(X\colon\Omega\fromto E\subseteq\realnumbers\) be a real-valued
    \discreterandomvariable. The \emph{expected value} of \(X\) exists if
    the series
    \[
        \sum_{\omega\in\Omega}X(\omega)\mathbb P(\{\omega\})
    \]
    is well-defined in \([-\infty,+\infty]\). In this case,
    \[
        \expectation[X]
        =
        \sum_{\omega\in\Omega}X(\omega)\mathbb P(\{\omega\}).
    \]
\end{snippetdefinition}

\begin{snippetdefinition}{finite-expected-value-definition}{Finite expected value}
    Let \((\Omega,\powerset(\Omega),\mathbb P)\) be a
    \countableprobabilityspace, and let
    \(X\colon\Omega\fromto E\subseteq\realnumbers\) be a real-valued
    \discreterandomvariable. The random variable \(X\) has a \emph{finite
    expected value} if
    \[
        \expectation[|X|]<+\infty.
    \]
    In this case \(\expectation[X]\in\realnumbers\).
\end{snippetdefinition}

\begin{snippetexample}{discrete-random-variable-infinite-expected-value-example}{A discrete random variable with infinite expected value}
    Let \(\Omega=\naturalnumbers^\exceptzero\), and let
    \((\Omega,\powerset(\Omega),\mathbb P)\) be the
    \countableprobabilityspace defined by
    \[
        \mathbb P(\{n\})=\frac{c}{n^2},
        \qquad
        c=\left(\sum_{n=1}^\infty\frac{1}{n^2}\right)^{-1}.
    \]
    Let \(X\colon\Omega\fromto\naturalnumbers^\exceptzero\) be
    \(X(n)=n\). Then
    \[
        \expectation[X]
        =
        \sum_{n=1}^\infty n\frac{c}{n^2}
        =
        c\sum_{n=1}^\infty\frac{1}{n}
        =
        +\infty.
    \]
\end{snippetexample}

\begin{snippettheorem}{expected-value-density-formula}{Expected value through the density}
    Let \((\Omega,\powerset(\Omega),\mathbb P)\) be a
    \countableprobabilityspace, and let
    \(X\colon\Omega\fromto E\subseteq\realnumbers\) be a real-valued
    \discreterandomvariable with \randomvariabledensity \(p_X\). Then
    \(X\) has a finite \expectedvalue if and only if
    \[
        \sum_{x\in E}|x|p_X(x)<+\infty.
    \]
    In this case,
    \[
        \expectation[X]=\sum_{x\in E}xp_X(x).
    \]
\end{snippettheorem}

\begin{snippetproof}{expected-value-density-formula-proof}{expected-value-density-formula}{Expected value through the density}
    For every \(x\in E\), let
    \[
        \Omega_x=\{\omega\in\Omega\suchthat X(\omega)=x\}.
    \]
    The \event[events] \(\Omega_x\), with \(x\in E\), are pairwise disjoint
    and their union is \(\Omega\). Moreover, \(p_X(x)=\mathbb P(\Omega_x)\).
    Hence, whenever the sums are finite,
    \begin{align*}
        \sum_{\omega\in\Omega}X(\omega)\mathbb P(\{\omega\})
        &=
        \sum_{x\in E}\sum_{\omega\in\Omega_x}
        X(\omega)\mathbb P(\{\omega\})
        \\
        &=
        \sum_{x\in E}x\sum_{\omega\in\Omega_x}\mathbb P(\{\omega\})
        =
        \sum_{x\in E}x\mathbb P(\Omega_x)
        =
        \sum_{x\in E}xp_X(x).
    \end{align*}
    Applying the same identity to \(|X|\) gives the finiteness criterion.
\end{snippetproof}

\begin{snippetproposition}{indicator-random-variable-expected-value}{Expected value of an indicator}
    Let \((\Omega,\powerset(\Omega),\mathbb P)\) be a
    \countableprobabilityspace, and let \(A\subseteq\Omega\). Then
    \[
        \expectation[\indicator_A]=\mathbb P(A).
    \]
\end{snippetproposition}

\begin{snippetproof}{indicator-random-variable-expected-value-proof}{indicator-random-variable-expected-value}{Expected value of an indicator}
    Since \(\indicator_A(\omega)=1\) exactly on \(A\),
    \[
        \expectation[\indicator_A]
        =
        \sum_{\omega\in A}\mathbb P(\{\omega\})
        =
        \mathbb P(A).
    \]
\end{snippetproof}

\section{Change of variable}

\begin{snippettheorem}{discrete-random-variable-change-of-variable-expectation}{Change-of-variable formula for discrete random variables}
    Let \((\Omega,\powerset(\Omega),\mathbb P)\) be a
    \countableprobabilityspace. Let
    \(X\colon\Omega\fromto E\subseteq\realnumbers^n\) be a
    \discreterandomvariable with \randomvariabledensity \(p_X\), and let
    \(g\colon E\fromto\realnumbers\) be a \function. Then \(g\circ X\) has a
    finite \expectedvalue if and only if
    \[
        \sum_{x\in E}|g(x)|p_X(x)<+\infty.
    \]
    In this case,
    \[
        \expectation[g(X)]
        =
        \sum_{x\in E}g(x)p_X(x).
    \]
\end{snippettheorem}

\begin{snippetproof}{discrete-random-variable-change-of-variable-expectation-proof}{discrete-random-variable-change-of-variable-expectation}{Change-of-variable formula for discrete random variables}
    For every \(x\in E\), let
    \[
        \Omega_x=\{\omega\in\Omega\suchthat X(\omega)=x\}.
    \]
    As in the density formula for the \expectedvalue,
    \begin{align*}
        \expectation[g(X)]
        &=
        \sum_{x\in E}\sum_{\omega\in\Omega_x}
        g(X(\omega))\mathbb P(\{\omega\})
        \\
        &=
        \sum_{x\in E}g(x)\mathbb P(\Omega_x)
        =
        \sum_{x\in E}g(x)p_X(x).
    \end{align*}
    Applying the same formula to \(|g\circ X|\) gives the finiteness
    criterion.
\end{snippetproof}

\begin{snippetcorollary}{joint-density-product-expectation-formula}{Expectation of a function of a discrete random vector}
    Let \((\Omega,\powerset(\Omega),\mathbb P)\) be a
    \countableprobabilityspace. For every \(i\in\{1,\ldots,n\}\), let
    \(X_i\colon\Omega\fromto E_i\subseteq\realnumbers\) be a
    \discreterandomvariable, and let \(p_{X_1,\ldots,X_n}\) be the
    \jointdensity of \((X_1,\ldots,X_n)\). Let
    \(g\colon\prod_{i=1}^n E_i\fromto\realnumbers\) be a \function. If
    \[
        \sum_{(x_1,\ldots,x_n)\in\prod_{i=1}^n E_i}
        |g(x_1,\ldots,x_n)|p_{X_1,\ldots,X_n}(x_1,\ldots,x_n)<+\infty,
    \]
    then
    \[
        \expectation[g(X_1,\ldots,X_n)]
        =
        \sum_{(x_1,\ldots,x_n)\in\prod_{i=1}^n E_i}
        g(x_1,\ldots,x_n)p_{X_1,\ldots,X_n}(x_1,\ldots,x_n).
    \]
\end{snippetcorollary}

\section{Properties}

\begin{snippetproposition}{basic-properties-discrete-expected-value}{Basic properties of expected value}
    Let \((\Omega,\powerset(\Omega),\mathbb P)\) be a
    \countableprobabilityspace, and let
    \(X,Y\colon\Omega\fromto\realnumbers\) be real-valued
    \discreterandomvariable[discrete random variables] with finite
    \expectedvalue[expected values]. Then:
    \begin{enumerate}
        \item if \(X\leq Y\) almost surely, then
            \(\expectation[X]\leq\expectation[Y]\);
        \item \(|\expectation[X]|\leq\expectation[|X|]\);
        \item for every \(a,b\in\realnumbers\),
            \[
                \expectation[aX+bY]
                =
                a\expectation[X]+b\expectation[Y];
            \]
        \item if \(X\geq0\) almost surely, then
            \[
                \expectation[X]=0
                \quad\Longleftrightarrow\quad
                X\asequal 0;
            \]
        \item if \(X\asequal c\) for some \(c\in\realnumbers\), then
            \[
                \expectation[X]=c.
            \]
    \end{enumerate}
\end{snippetproposition}

\begin{snippetproof}{basic-properties-discrete-expected-value-proof}{basic-properties-discrete-expected-value}{Basic properties of expected value}
    Each statement follows from the defining series. Linearity follows by
    distributing the absolutely convergent sums. Monotonicity follows because
    the summands of \(Y-X\) are nonnegative outside a null \event. The
    inequality \(|\expectation[X]|\leq\expectation[|X|]\) is the triangle
    inequality for the series. If \(X\geq0\) and \(\expectation[X]=0\), all
    nonnegative summands \(X(\omega)\mathbb P(\{\omega\})\) vanish, so
    \(X=0\) outside a null \event. The constant case is immediate.
\end{snippetproof}

\begin{snippetproposition}{equality-in-law-preserves-discrete-expectation}{Equality in law preserves expected value}
    Let
    \((\Omega_1,\powerset(\Omega_1),\mathbb P_1)\) and
    \((\Omega_2,\powerset(\Omega_2),\mathbb P_2)\) be
    \countableprobabilityspace[countable probability spaces]. Let
    \(X\colon\Omega_1\fromto E\subseteq\realnumbers\) and
    \(Y\colon\Omega_2\fromto E\) be real-valued
    \discreterandomvariable[discrete random variables]. If \(X\eqdist Y\),
    then \(X\) has a finite \expectedvalue if and only if \(Y\) has a finite
    \expectedvalue, and in that case
    \[
        \expectation[X]=\expectation[Y].
    \]
\end{snippetproposition}

\begin{snippetproof}{equality-in-law-preserves-discrete-expectation-proof}{equality-in-law-preserves-discrete-expectation}{Equality in law preserves expected value}
    If \(X\eqdist Y\), then \(p_X(x)=p_Y(x)\) for every \(x\in E\). The
    density formula for the \expectedvalue gives both the finiteness
    equivalence and the equality of the two expected values.
\end{snippetproof}

\end{document}
